\section{Endpoint deletion and the full-kernel bridge}
\label{sec:endpoint}

The zero mode \eqref{eq:exact-zero-mode} is initially a continuous integral
with the discrete restrictions $l,n\geq1$ retained by $U$.  We first delete
the artificial regions $0<l<1$ and $0<n<1$ on a genuine convergence patch.
We then construct the exact continuation needed to reach all small shifts.

\subsection{Fourier decay and the endpoint lemma}

\begin{lemma}[Fourier--conductor bound]\label{lem:fourier-conductor}
Uniformly for $x,y>0$, $x\neq y$, and every $A_0,J>0$,
\begin{align}
 &\left|\int_{\R}w(t)(x/y)^{it}
 \V^\pm\!\left(\frac{y}{kmY},\frac{mx}{h},t\right)\dd t\right|
 \notag\\
 &\qquad\ll_{A_0,J}T(\log T)^{C_M}
 \left(1+\frac{y}{kmY}\right)^{-A_0}\notag\\
 &\hspace{35mm}\times
 \left(1+\frac{mx}{hT}\right)^{-A_0}
 \left(1+\Delta|\log(x/y)|\right)^{-J}.
 \label{eq:fourier-conductor}
\end{align}
\end{lemma}

\begin{proof}
Integrate by parts $J$ times in $t$ and use
\eqref{eq:w-derivatives} and \eqref{eq:double-weight-bound}.  Every
derivative costs at most $\Delta^{-1}$, and the support of $w$ has length
$O(T)$.
\end{proof}

\begin{proposition}[Endpoint deletion on the good half-space]
\label{prop:endpoint-deletion}
Assume \eqref{eq:main-range}.  For every $D>0$, after summing over
$h\leq H$, $k\leq K$ with the outside coefficient weights, one has
\begin{equation}\label{eq:endpoint-deletion}
 \mathcal O^\pm_{0,Y}-\widetilde{\mathcal O}^\pm_{0,Y}
 =O_D(T^{-D})
\end{equation}
when
\begin{equation}\label{eq:endpoint-halfspace}
 \Re(\beta-\gamma)\geq0,\qquad
 |\Re\alpha|+|\Re\beta|+|\Re\gamma|\leq C/\log T.
\end{equation}
The source expressions retain the complete contours in
\eqref{eq:coupled-weight}.
\end{proposition}

\begin{proof}
It suffices to treat
\begin{equation}\label{eq:endpoint-regions}
             \mathcal E_l=\{y<km\},\qquad
             \mathcal E_n=\{x<h\}.
\end{equation}
Put $\delta_0=\eta/4$ and $\eps_0=\eta/8$.  On $\mathcal E_l$, suppose
simultaneously that
\begin{equation}\label{eq:slow-small-assumption}
 |\log(x/y)|\leq T^{-1+\delta_0},\qquad
 \frac{mx}{hT}\leq T^{\eps_0}.
\end{equation}
Then $x\asymp y$, and, with $r=d\rho$, $h=dq$, $km=dv$,
\begin{equation}\label{eq:scale-rho}
 y\gg d|\rho|T^{1-\delta_0},\qquad
 v\gg|\rho|T^{1-\delta_0},\qquad
 m|\rho|\ll qT^{\delta_0+\eps_0}.
\end{equation}
Since $v\leq km$ and $q\leq h\leq H$, this forces
\begin{equation}\label{eq:scale-contradiction}
 HK\gg T^{1-2\delta_0-\eps_0}=T^{1-5\eta/8}.
\end{equation}
But \eqref{eq:main-range} and $H\geq1$ give $HK\leq T^{1-\eta}$, a
contradiction for large $T$.  Every point of $\mathcal E_l$ therefore has
either Fourier saving $|\log(x/y)|>T^{-1+\delta_0}$ or conductor saving
$mx/(hT)>T^{\eps_0}$.  The orders $A_0,J$ in
\eqref{eq:fourier-conductor} absorb every remaining polynomial factor.

On $\mathcal E_n$, $0<x<h$ and $r=d\rho\neq0$ imply
\begin{equation}\label{eq:n-endpoint-phase}
 |\log(x/y)|\gg\min(1,|r|/h)\gg 1/H.
\end{equation}
The main range gives $H\leq T^{2(1-\eta)/3}$, and hence
\begin{equation}\label{eq:n-endpoint-saving}
 \Delta|\log(x/y)|\gg T^{(1+2\eta)/3}/\log T.
\end{equation}
Repeated integration by parts again gives an arbitrary power saving, once
the endpoint mass is shown to be only polynomial.

We record the singular-layer summation because it is where the gcd density is
needed.  For fixed $m$ and a dyadic $l$-scale, the values
\[
                 n=\frac{kml+d\rho}{h}
\]
form a mesh of spacing $d/h$, exactly matching
\eqref{eq:poisson-density}.  A boundary layer meeting $n=0$ has
$kml\asymp d|\rho|$.  On writing $n=(d/h)x$ and $x=|\rho|z$, the square-root
Jacobian is
\begin{equation}\label{eq:square-root-jacobian}
 |\rho|^{-1/2}\int_0^{O(|\rho|)}x^{-1/2}(\cdots)\dd x
 =\int_0^{O(1)}z^{-1/2}(\cdots)\dd z.
\end{equation}
Thus no unrecorded sum over $\rho$ and no lost $|\rho|^{-1/2}$ occurs.

For the potentially worst case $r<0$, $x<h$, partial summation in
$r=d\rho$ and then in $m$ gives, up to $T^\eps$,
\begin{align}
 \sum_{\rho\geq1}y^{-1/2-\Re\alpha}
 \left(1+\frac{y}{kmY}\right)^{-A_0}
 &\ll d^{-1}(kmY)^{1/2-\Re\alpha},\label{eq:rho-endpoint-sum}\\
 \sum_{m\geq1}m^{-1/2-\Re\beta}
 \left(1+\frac{mx}{hT}\right)^{-A_0}
 &\ll(hT/x)^{1/2-\Re\beta}.
 \label{eq:m-endpoint-sum}
\end{align}
After multiplication by $x^{-1/2-\Re\gamma}$, the lower-end behavior is
\begin{equation}\label{eq:endpoint-integrability}
                          x^{-1+\Re(\beta-\gamma)}.
\end{equation}
It is integrable in \eqref{eq:endpoint-halfspace}; on the boundary, the last
factor in \eqref{eq:fourier-conductor} makes the logarithmic endpoint
integrable for $J>1$.  All remaining sums have fixed polynomial size and are
absorbed by \eqref{eq:n-endpoint-saving}.  This proves
\eqref{eq:endpoint-deletion}.
\end{proof}

The ordinary dominated-convergence proof stops at
$\Re(\beta-\gamma)=0$: its far dyadic majorant has residual size
\begin{equation}\label{eq:far-tail-obstruction}
 M_1^{\Re(\gamma-\beta)}
 \{1+\log(M_1/T)\}^{-J}\qquad(M_1\gg T).
\end{equation}
We therefore continue only the complete assembled error, rather than a
single endpoint-deleted branch.

\subsection{Source normality with the full kernel}

In this subsection use external shifts $A,B,C_1$ in place of
$\alpha,\beta,\gamma$ in the kernel.  The continuation requires only one
nonempty patch; for definiteness take $\kappa_0=1$, $c_1=5$, and put
\begin{equation}\label{eq:good-patch}
 \Omega_g=\left\{
 \Re B>\frac{\kappa_0}{\log T},\quad
 \Re C_1<-\frac{\kappa_0}{\log T},\quad
 |A|+|B|+|C_1|<\frac{c_1}{\log T}\right\}.
\end{equation}
Fix $B,C_1$ there and let $A$ vary in the connected strip
\begin{equation}\label{eq:external-A-strip}
                       \Re A<\frac12-\frac{c_u}{2}.
\end{equation}

\begin{lemma}[Full-kernel source normality]\label{lem:source-normality}
For each of the four pairs
\[(\text{direct or dual AFE sign},\ \operatorname{sgn}\rho),\]
the endpoint-deleted expressions, with the finite regulators introduced
below, converge locally uniformly in the strip
\eqref{eq:external-A-strip}, uniformly for $(B,C_1)$ on compact subsets of
\eqref{eq:good-patch}.  Their limits are holomorphic in $A$, as are any fixed
number of $A$-derivatives.
\end{lemma}

\begin{proof}
Keep the complete kernel $\V(l/Y,mn,t)$ and do not freeze its $s$ variable.
For the exact-cutoff source, $l,n\geq1/2$.  If $F(n)$ is the nonnegative
product-tail majorant from \Cref{lem:double-weight}, then on either half-mesh
of spacing $\delta=d/h$,
\begin{equation}\label{eq:mesh-majorant}
 \delta\sum_{\rho:n\geq1/2}F(n)
 \ll\int_{1/2}^{\infty}F(v)\dd v
     +\delta\sup_{v\geq1/2}F(v).
\end{equation}
This gives, before the finite outer sums, a fixed logarithmic power times
\begin{multline}
 T\int_{1/2}^{\infty}l^{-1/2-a_0}(1+l/Y)^{-P}\dd l
 \sum_{m\geq1}m^{-1/2-b_0}\\
 {}\times\left\{
 \int_{1/2}^{\infty}n^{-1/2-c_0}(1+mn/T)^{-Q}\dd n
 +(1+m/T)^{-Q}\right\}.
 \label{eq:source-product-majorant}
\end{multline}
where $a_0=\Re A$, $b_0=\Re B$, and $c_0=\Re C_1$.  On compact subsets of
\eqref{eq:external-A-strip}, this is polynomial in $T$ and locally uniform;
the residual $m$-sum is $O(\log T)$ on the good patch.  The same is true
after the dual substitution $(B,C_1)\mapsto(-C_1,-B)$.

For completeness, consider the worst $n=0$ discrepancy with
$r=-d\rho$, $z=x/(d\rho)\leq1$, and restore the outside coefficient weight.
Its full algebraic prefactor is
\begin{multline}
 |a_hb_k|d^{1-a_0-c_0}h^{-1+c_0}k^{-1+a_0}
 m^{-1+a_0-b_0}\rho^{-a_0-c_0}\\
 {}\times\int_0^1z^{-1/2-c_0}
 \left(1+\frac{md\rho z}{hT}\right)^{-Q}
 \left(1+\frac{d\rho}{kmY}\right)^{-P}
 (1+|\log z|)^{-J}\dd z.
 \label{eq:worst-layer-prefactor}
\end{multline}
For $X=md\rho/(hT)$,
\begin{equation}\label{eq:z-integral-bound}
 \int_0^1z^{-1/2-c_0}(1+Xz)^{-Q}(1+|\log z|)^{-J}\dd z
 \ll(1+X)^{-1/2+c_0}\{1+\log(2+X)\}^{-J}.
\end{equation}
In the range $X>1$, summing $m$ and $\rho$ gives at most
\begin{equation}\label{eq:mrho-bound}
 \left(\frac{hT}{d}\right)^{1/2-c_0}
 \left(\frac{kY}{d}\right)^{1/2-a_0}
 \sum_{m\geq1}m^{-1-b_0+c_0}.
\end{equation}
The last sum is $O_{\kappa_0}(\log T)$.  Multiplication by the prefactor in
\eqref{eq:worst-layer-prefactor} cancels every power of $d$ and leaves
\begin{equation}\label{eq:d-cancellation}
 |a_hb_k|(hk)^{-1/2}T^{1/2-c_0}Y^{1/2-a_0}
 \sum_{m\geq1}m^{-1-b_0+c_0}.
\end{equation}
After the $h,k$ sums and the Fourier prefactor, the $X>1$ discrepancy is
\begin{equation}\label{eq:full-worst-layer-bound}
 \ll T\Delta^{-J}T^{1/2-c_0}Y^{1/2-a_0}
 (HK)^{1/2}T^\eps\log T,
\end{equation}
and is $O_N(T^{-N})$ after increasing $J$.  This exact cancellation is
uniform over all gcd blocks.  The range $X\leq1$
has $m\rho\leq hT/d$ and therefore only polynomial mass on compact external
$A$-strips.  For $z>1$, $x<h=dq$ gives
\begin{equation}\label{eq:z-large-phase}
 \left|\log\frac{x}{x+d\rho}\right|
 =\log(1+z^{-1})\gg q^{-1}\geq H^{-1}.
\end{equation}
In both cases an arbitrarily high power in the Fourier bound absorbs the
polynomial mass.  On far $m$-scales, the residual exponent is
\begin{equation}\label{eq:far-m-exponent}
 a_0-\Re(B-C_1)-\frac12
 \leq-\frac{c_u}{2}-\Re(B-C_1),
\end{equation}
so the dyadic sum is uniform on the good patch.  The estimates remain valid
after $A$-derivatives, which only insert logarithms.  This proves local
normality and holomorphy.
\end{proof}

\subsection{Finite regulators and genuine Tonelli domains}

\begin{proposition}[Four-branch regulator--Tonelli bridge]
\label{prop:four-branch-regulator-tonelli}
Fix $(B,C_1)$ in a compact subset of \eqref{eq:good-patch}, keep
$\Re u=c_u$, and put $\lambda=A+u$.  Define
\begin{align}
 z_{\rm d}&=\lambda+C_1+s,
 &\nu_{\rm d}&=1-\lambda+B+s,
 \label{eq:four-branch-direct-vars}\\
 z_{\rm q}&=\lambda-B+s,
 &\nu_{\rm q}&=1-\lambda-C_1+s.
 \label{eq:four-branch-dual-vars}
\end{align}
Here ${\rm d}$ denotes the direct AFE piece and ${\rm q}$ the dual piece.
In each of the four branches impose
\begin{equation}\label{eq:finite-regulators}
 \begin{gathered}
 m,|\rho|\leq R,\qquad R^{-1}\leq l,n\leq R,\\
 e^{-\kappa(m+|\rho|+l+n)},\qquad
 R\geq2,\quad0<\kappa\leq1.
 \end{gathered}
\end{equation}
Then the source limits obtained in the order
\begin{equation}\label{eq:four-branch-order}
 V\longrightarrow\infty,\qquad R\longrightarrow\infty,
 \qquad\kappa\longrightarrow0^+
\end{equation}
exist locally uniformly in the connected strip
\eqref{eq:external-A-strip} and are holomorphic in the original shift $A$.
The corresponding destination limits are absolutely convergent in the
following four nonempty domains.  The two negative branches use the common
corridor
\begin{equation}\label{eq:lambda-corridor}
            \frac12<\Re\lambda
            <\frac12+\Re(B-C_1).
\end{equation}
For the direct positive and negative branches, respectively, the remaining
conditions are
\begin{equation}\label{eq:positive-direct-tonelli}
 \begin{aligned}
 \Re\lambda&<\frac12,\\
 \Re s&>\max\{1-\Re(\lambda+C_1),\Re(\lambda-B)\}.
 \end{aligned}
\end{equation}
\begin{equation}\label{eq:direct-tonelli}
 \max\{1-\Re(\lambda+C_1),\Re(\lambda-B)\}
 <\Re s<\frac12-\Re C_1.
\end{equation}
For the dual positive and negative branches, respectively, they are
\begin{equation}\label{eq:positive-dual-tonelli}
 \begin{aligned}
 \Re\lambda&<\frac12,\\
 \Re s&>\max\{1-\Re(\lambda-B),\Re(\lambda+C_1)\}.
 \end{aligned}
\end{equation}
\begin{equation}\label{eq:dual-tonelli}
 \max\{1-\Re(\lambda-B),\Re(\lambda+C_1)\}
 <\Re s<\frac12+\Re B.
\end{equation}
For each branch, the regulated source and destination agree exactly.  After
the limits in \eqref{eq:four-branch-order}, this equality extends from the
relevant Tonelli open set to \eqref{eq:external-A-strip} by the
one-variable identity theorem in $A$, with $B,C_1$ fixed.  No continuation
in the separate symbols $A+u$ and $s$ is used.
\end{proposition}

\begin{proof}
Write $\mathfrak r=d|\rho|>0$.  After the endpoint cutoffs have been deleted,
the four continuous integrals determining branchwise convergence are as
follows.  They are displayed separately because their zero endpoints differ.

For the direct, positive-$\rho$ branch, $x=y+\mathfrak r$ and
\begin{align}
 I_{{\rm d},+}(\mathfrak r)
 &:=\int_0^\infty y^{-1/2-\lambda-it}
       (y+\mathfrak r)^{-1/2-C_1+it-s}\dd y\notag\\
 &=\mathfrak r^{-z_{\rm d}}
   \frac{\Gamma(\frac12-\lambda-it)\Gamma(z_{\rm d})}
        {\Gamma(\frac12+C_1-it+s)}.
 \label{eq:four-beta-dp}
\end{align}
At $y=0$ this requires $\Re\lambda<1/2$; at infinity it first requires
$\Re z_{\rm d}>0$.  The subsequent $\rho$-sum strengthens the latter to
$\Re z_{\rm d}>1$.  On that half-plane the gcd exponent
$1-z_{\rm d}$ has negative real part, and the remaining $m$-sum is
absolutely convergent when $\Re\nu_{\rm d}>1$.  These conditions are
exactly \eqref{eq:positive-direct-tonelli}.

For the direct, negative-$\rho$ branch, put $y=x+\mathfrak r$.  Then
\begin{align}
 I_{{\rm d},-}(\mathfrak r)
 &:=\int_0^\infty(x+\mathfrak r)^{-1/2-\lambda-it}
       x^{-1/2-C_1+it-s}\dd x\notag\\
 &=\mathfrak r^{-z_{\rm d}}
   \frac{\Gamma(\frac12-C_1+it-s)\Gamma(z_{\rm d})}
        {\Gamma(\frac12+\lambda+it)}.
 \label{eq:four-beta-dm}
\end{align}
The endpoint $x=0$ gives $\Re(C_1+s)<1/2$, while infinity, the
$\rho$-sum, and the $m$-sum give respectively
$\Re z_{\rm d}>0$, $\Re z_{\rm d}>1$, and
$\Re\nu_{\rm d}>1$.  Hence
\[
 \max\{1-\Re(\lambda+C_1),\Re(\lambda-B)\}
 <\Re s<\frac12-\Re C_1.
\]
The first lower bound is below the upper bound precisely when
$\Re\lambda>1/2$; the second is below it precisely when
$\Re\lambda<1/2+\Re(B-C_1)$.  This proves
\eqref{eq:direct-tonelli}, including the necessity of its corridor.

The dual AFE piece makes the substitution
$(B,C_1)\mapsto(-C_1,-B)$ and has the fixed exterior factor
$X_{B,C_1}(t)$.  Its positive-$\rho$ integral is, explicitly,
\begin{align}
 I_{{\rm q},+}(\mathfrak r)
 &:=\int_0^\infty y^{-1/2-\lambda-it}
       (y+\mathfrak r)^{-1/2+B+it-s}\dd y\notag\\
 &=\mathfrak r^{-z_{\rm q}}
   \frac{\Gamma(\frac12-\lambda-it)\Gamma(z_{\rm q})}
        {\Gamma(\frac12-B-it+s)}.
 \label{eq:four-beta-qp}
\end{align}
The $y=0$ condition is $\Re\lambda<1/2$; the infinity, $\rho$-, and
$m$-conditions are respectively
$\Re z_{\rm q}>0$, $\Re z_{\rm q}>1$, and
$\Re\nu_{\rm q}>1$.  They give
\eqref{eq:positive-dual-tonelli}.

Finally, the dual, negative-$\rho$ branch is
\begin{align}
 I_{{\rm q},-}(\mathfrak r)
 &:=\int_0^\infty(x+\mathfrak r)^{-1/2-\lambda-it}
       x^{-1/2+B+it-s}\dd x\notag\\
 &=\mathfrak r^{-z_{\rm q}}
   \frac{\Gamma(\frac12+B+it-s)\Gamma(z_{\rm q})}
        {\Gamma(\frac12+\lambda+it)}.
 \label{eq:four-beta-qm}
\end{align}
Here $x=0$ requires $\Re s<1/2+\Re B$; infinity and the two sums require
$\Re z_{\rm q}>1$ and $\Re\nu_{\rm q}>1$.  Therefore
\[
 \max\{1-\Re(\lambda-B),\Re(\lambda+C_1)\}
 <\Re s<\frac12+\Re B.
\]
The two comparisons with the upper endpoint give separately
$\Re\lambda>1/2$ and
$\Re\lambda<1/2+\Re(B-C_1)$, proving
\eqref{eq:dual-tonelli}.  Since $\Re(B-C_1)>0$ on the good patch, both
negative intervals are nonempty.

We next justify the contour operation before evaluating any beta integral.
Keep $R,\kappa$ fixed and truncate the complete $s$- and $u$-lines at a
common height $V$.  At this finite-regulator stage the possible
singularities of the unreduced source Mellin integrand are as follows.
\begin{enumerate}[label=\textup{(\alph*)},leftmargin=8mm]
\item The pole $s=0$ comes from $1/s$.  Every source and destination
$s$-line just used has positive real part, so it is not crossed.
\item The poles $u=0,-1,-2,\ldots$ come from $\Gamma(u)$.  The line
$\Re u=c_u>0$ is not moved, so none is crossed.
\item The numerator-gamma poles of $g_{B,C_1}$ are
\[
 s=-\frac12-B-it-2j,\qquad s=-\frac12-C_1+it-2j
 \quad(j\geq0),
\]
and those of $g_{-C_1,-B}$ are
\[
 s=-\frac12+C_1-it-2j,\qquad s=-\frac12+B+it-2j
 \quad(j\geq0).
\]
For large $T$ their real parts lie to the left of every positive line used
above.  Denominator gamma factors give zeros, not poles.
\item The kernel $\mathcal G(s,u)$ is entire in $(s,u)$; its displayed
denominators depend only on the external shifts and stay bounded away from
zero on the good patch.  Because $m,|\rho|$ are finite at this stage, there
is no arithmetic zeta function and hence no $z=1$ or $\nu=1$ pole.
\end{enumerate}
Thus the finite rectangles deform without a residue.  On a horizontal
$s$-edge, $s=\sigma+iV$,
\[
                         |e^{s^2}|=e^{\sigma^2-V^2},
\]
while the gamma ratios and remaining factors have at most exponential times
polynomial growth in $V$.  On a horizontal $u$-edge, the integrand is still
the unreduced source integral, and Stirling gives
\[
 |\Gamma(c_u+iV)|\ll(1+V)^{c_u-1/2}e^{-\pi V/2}.
\]
All other $u$-dependence is polynomial at fixed $R,\kappa$.  Hence all
horizontal sides and vertical tails tend to zero as $V\to\infty$, uniformly
for the fixed regulators.  The complete vertical contours are thereby
restored before a variable regulator is removed.

At a Tonelli destination, first let $R\to\infty$ with $\kappa>0$ fixed and
then let $\kappa\to0^+$.  Absolute convergence in
\eqref{eq:positive-direct-tonelli}--\eqref{eq:dual-tonelli} justifies both
limits and, only after the second, the beta evaluations
\eqref{eq:four-beta-dp}--\eqref{eq:four-beta-qm}.  In particular, no beta
formula is asserted with the Abel factor still present.  At the source the
same order is justified by \Cref{lem:source-normality}: the complete coupled
kernel, exact mesh density, and product-tail majorant give bounds locally
uniform in $A$ and independent of $R,\kappa$.  This proves the locally
uniform holomorphic source limit in \eqref{eq:external-A-strip}.

For completeness, singularities appearing only after the variable
regulators have been removed must be recorded separately.  The direct
branches produce $\zeta(z_{\rm d})$ and a gcd--$m$ factor with polar part
$\zeta(\nu_{\rm d})$; the dual branches similarly produce
$\zeta(z_{\rm q})$ and $\zeta(\nu_{\rm q})$.
\begin{enumerate}[label=\textup{(\roman*)},leftmargin=8mm]
\item The direct pole $z_{\rm d}=1$ is cancelled by the kernel zero
$s=1-\lambda-C_1$, and the dual pole $z_{\rm q}=1$ by the zero
$s=1-\lambda+B$.
\item After applying the zeta functional equation, the direct hyperplanes
$z_{\rm d}=0$, $\nu_{\rm d}=1$, and the corresponding dual hyperplanes
$z_{\rm q}=0$, $\nu_{\rm q}=1$ remain.  They are not discarded or crossed
branchwise: these are the small arithmetic divisors cancelled only after
the four branches and both diagonals are assembled.
\item Before the positive- and negative-$\rho$ pieces are combined, the beta
factors may have the following gamma poles, where $j\geq0$:
\begin{align*}
 ({\rm d},+):&\quad z_{\rm d}=-j,\quad
   \lambda=\frac12+j-it;\\
 ({\rm d},-):&\quad z_{\rm d}=-j,\quad
   s=\frac12-C_1+it+j;\\
 ({\rm q},+):&\quad z_{\rm q}=-j,\quad
   \lambda=\frac12+j-it;\\
 ({\rm q},-):&\quad z_{\rm q}=-j,\quad
   s=\frac12+B+it+j.
\end{align*}
They lie outside the four Tonelli domains.  After combining the two
$\rho$-signs, \Cref{lem:gamma-collapse} cancels every apparent near-real
archimedean pole.  The remaining direct poles are
\[
 s=n-\frac12-C_1+it\ (n\geq1),\qquad
 s=-\frac12-B-it-2j,\qquad s=-\frac12-C_1+it-2j,
\]
and the remaining dual poles are
\[
 s=n-\frac12+B+it\ (n\geq1),\qquad
 s=-\frac12+C_1-it-2j,\qquad s=-\frac12+B+it-2j.
\]
All have imaginary height comparable with $T$.  If a later fixed-line
deformation crosses one, $e^{s^2}$ makes its residue smaller than every
power of $T$.  Reciprocal gamma factors again give zeros only.
\end{enumerate}
The poles $s=0$ and those of $\Gamma(u)$ have not moved during this bridge;
they are treated in the subsequent fixed-$u$ assembly and final $u$-move.

It remains to identify the selected continuation.  The positive branches
agree with their destinations on the nonempty open set
\[
 \mathscr U_+=\{A:\Re A+c_u<1/2\}
                \cap\{\Re A<1/2-c_u/2\},
\]
whereas the negative branches agree on
\[
 \mathscr U_-=\left\{A:\frac12<\Re A+c_u
       <\frac12+\Re(B-C_1)\right\}
                \cap\{\Re A<1/2-c_u/2\}.
\]
The latter is nonempty because $\Re(B-C_1)>0$; for large $T$ its width is
$O(1/\log T)<c_u/2$, so it lies in the connected strip
\eqref{eq:external-A-strip}.  Apply the one-variable identity theorem
separately to each complete, already integrated branch as a function of
$A$.  This extends an exact equality to the target $A$; it does not
propagate the endpoint-deletion $O(T^{-N})$ error and does not treat
$\lambda$ and $s$ as independent continuation variables.
\end{proof}
